% !TEX root = main.tex
\section{Introduction}

Most modern control and instrumentation systems are digital: sensors produce continuous signals,
but the decisions are taken by processors that work with discrete samples and finite-length
binary numbers~\cite{tocci2011}. Converting a continuous physical process into this digital domain
introduces sampling, quantization and finite resolution in both the measurement and the actuation,
and the control law itself becomes a difference equation executed at fixed time
intervals~\cite{franklin1998,oppenheim2010}.

This paper uses a compact thermal chamber as a case study of a complete digital system. The chamber
is an acrylic box heated by three 12\,V/5\,W automotive lamps and controlled by an ESP32
microcontroller~\cite{esp32ds}. Instead of focusing on the hardware, the analysis follows the
signal through each digital stage, from the analog-to-digital conversion to the PWM output, and
describes the finite state machine that organizes the operation of the system.
